\documentclass[../../main.tex]{subfiles}% 注意这里的文件路径不能用 ./main.tex ,否则用latexmk编译子文件会报错
\graphicspath{{\subfix{./image/}}} % 指定图片目录,后续可以直接使用图片文件名
% 注意这里的文件路径不能用 ../../image/ ,否则用latexmk编译子文件会报错

% 例如:
% \begin{figure}[H]
% \centering
% \includegraphics[scale=0.3]{图.png}
% \caption{}
% \label{figure:图}
% \end{figure}
% 注意:上述\label{}一定要放在\caption{}之后,否则引用图片序号会只会显示??.

\begin{document}

\section{线性算子的谱}

\begin{definition}
设$\mathscr{X}$为$B$空间，$A:D(A)\subseteq\mathscr{X}\to\mathscr{X}$为闭线性算子。
\begin{enumerate}[(1)]
\item 若存在$x_0\in D(A)\setminus\{\theta\}$，使得$Ax_0=\lambda x_0$，则称$\lambda\in\mathbb{C}$为$A$的\textbf{特征值}，$x_0$为$\lambda$对应的\textbf{特征元}。
\item 称集合
\begin{align*}
\rho(A)\triangleq\{\lambda\in\mathbb{C}\mid (\lambda I-A)^{-1}\text{存在且}(\lambda I-A)^{-1}\in\mathscr{L}(\mathscr{X})\}
\end{align*}
为$A$的\textbf{预解集}，$\rho(A)$中的元素称为$A$的\textbf{正则值}。
\item 称$\sigma(A)\triangleq\mathbb{C}\setminus\rho(A)$为$A$的\textbf{谱集}，$\sigma(A)$中的元素称为\textbf{谱点}。进一步划分：
\begin{enumerate}[(i)]
\item \textbf{点谱}$\sigma_p(A)$：$\lambda$为特征值，即$(\lambda I-A)^{-1}$不存在；
\item \textbf{连续谱}$\sigma_c(A)$：$(\lambda I-A)^{-1}$存在，$R(\lambda I-A)\neq\mathscr{X}$但$\overline{R(\lambda I-A)}=\mathscr{X}$；
\item \textbf{剩余谱}$\sigma_r(A)$：$(\lambda I-A)^{-1}$存在，$\overline{R(\lambda I-A)}\neq\mathscr{X}$。
\end{enumerate}
\end{enumerate}
\end{definition}
\begin{remark}
显然,当$\text{dim}\,\mathscr{X}<\infty$时,只有两种可能:
\begin{enumerate}
\item $\lambda$是$A$的特征值；

\item $(\lambda I-A)^{-1}$作为矩阵存在,即$(\lambda I-A)^{-1}\in \mathscr{L}(\mathscr{X})$.
\end{enumerate}
当$\text{dim}\,\mathscr{X}=\infty$时,就可能会出现连续谱和点谱.
\end{remark}

\begin{proposition}
设$\mathscr{X}$为$B$空间，$A:D(A)\subseteq\mathscr{X}\to\mathscr{X}$为闭线性算子,则$\sigma_p(A),\sigma_c(A),\sigma_r(A)$两两无交,并且满足
$$\sigma(A)=\sigma_p(A)\cup\sigma_c(A)\cup\sigma_r(A).$$
\end{proposition}
\begin{proof}

由定义知$\sigma_p(A),\sigma_c(A),\sigma_r(A)$两两无交是显然的.

设 $\lambda \in \sigma(A)$，则 $\lambda \notin \rho(A)$，即 $(\lambda I - A)^{-1} \notin \mathscr{L}(\mathscr{X})$。分两种情况：
\begin{enumerate}[(1)]
\item $(\lambda I - A)^{-1}$ 不存在。则直接由定义 $\lambda \in \sigma_p(A)$。

\item $(\lambda I - A)^{-1}$ 存在。记 $S = (\lambda I - A)^{-1}$，其定义域为 $D(S) = R(\lambda I - A)$。由于 $\lambda \notin \rho(A)$，$S$ 不能属于 $\mathscr{L}(\mathscr{X})$。因为 $A$ 是闭算子，$\lambda I - A$ 也是闭算子，从而其逆 $S$ 是闭算子。若 $R(\lambda I - A) = \mathscr{X}$，则由$\mathscr{X}$完备和\rrefpro{proposition:完备度量空间的子空间等价于闭集}{proposition:完备度量空间的子空间等价于闭集-2}知$D(S)=R(\lambda I-A)=\mathscr{X}$是闭的.于是由\hyperref[theorem:泛函分析--闭图像定理]{闭图像定理}知$S$连续，再由\refpro{proposition:线性算子有界当且仅当连续}知$S \in \mathscr{L}(\mathscr{X})$，矛盾!故必有 $R(\lambda I - A) \neq \mathscr{X}$。此时根据值域是否稠密划分：
\begin{enumerate}[(i)]
\item 若 $\overline{R(\lambda I - A)} = \mathscr{X}$，则 $\lambda \in \sigma_c(A)$；

\item 若 $\overline{R(\lambda I - A)} \neq \mathscr{X}$，则 $\lambda \in \sigma_r(A)$。
\end{enumerate}
故$\lambda \in \sigma_p(A) \cup \sigma_c(A) \cup \sigma_r(A)$。因此 $\sigma(A) \subseteq \sigma_p(A) \cup \sigma_c(A) \cup \sigma_r(A)$。
\end{enumerate}

设$\lambda \in \sigma_p(A) \cup \sigma_c(A) \cup \sigma_r(A)$,则
\begin{enumerate}[(1)]
\item 若 $\lambda \in \sigma_p(A)$，则逆不存在，故 $\lambda \notin \rho(A)$，即 $\lambda \in \sigma(A)$。

\item 若 $\lambda \in \sigma_c(A) \cup \sigma_r(A)$，则逆存在但 $D(S)=R(\lambda I - A) \neq \mathscr{X}$。从而 $(\lambda I - A)^{-1} \notin \mathscr{L}(\mathscr{X})$，故 $\lambda \notin \rho(A)$，即 $\lambda \in \sigma(A)$。
\end{enumerate}
因此$\sigma_p(A) \cup \sigma_c(A) \cup \sigma_r(A) \subseteq \sigma(A)$。

综上可知$\sigma(A) = \sigma_p(A) \cup \sigma_c(A) \cup \sigma_r(A)$.

\end{proof}

\begin{example}
设$\mathscr{X}=L^2[0,1]$，算子$A:u(t)\mapsto -\frac{\mathrm{d}^2}{\mathrm{d}t^2}u(t)$，定义
\begin{align*}
D(A)=\{u\in\mathscr{X}\mid Au\in\mathscr{X}\},
\end{align*}
则$A:D(A)\to\mathscr{X}$为闭线性算子，且
\begin{align*}
\sigma(A)=\sigma_p(A)=\{(2n\pi)^2\mid n=0,1,2,\cdots\}.
\end{align*}
\begin{proof}

一方面，对任意$n=0,1,2,\cdots$，有
\begin{align*}
-\frac{\mathrm{d}^2}{\mathrm{d}t^2}\begin{pmatrix}\sin2n\pi t\\\cos2n\pi t\end{pmatrix}=(2n\pi)^2\begin{pmatrix}\sin2n\pi t\\\cos2n\pi t\end{pmatrix},
\end{align*}
故$(2n\pi)^2\in\sigma_p(A)$。

另一方面，若$\lambda\neq(2n\pi)^2$，对任意$f\in L^2[0,1]$，方程
\begin{align*}
\left(-\frac{\mathrm{d}^2}{\mathrm{d}t^2}-\lambda\right)u(t)=f(t)
\end{align*}
有唯一解
\begin{align*}
u(t)=\sum\limits_{n=-\infty}^{\infty}\frac{C_n}{(2n\pi)^2-\lambda}\mathrm{e}^{2\pi\mathrm{i}nt},
\end{align*}
其中$C_n=\int_{0}^{1}f(t)\mathrm{e}^{-2\pi\mathrm{i}nt}\mathrm{d}t,\ n\in\mathbb{Z}$。
易验证$u\in D(A)$，且
\begin{align*}
\|u\|^2=\sum\limits_{n=-\infty}^{\infty}\frac{|C_n|^2}{|(2n\pi)^2-\lambda|^2}\leqslant M_{\lambda}^2\sum\limits_{n=-\infty}^{\infty}|C_n|^2=M_{\lambda}^2\|f\|^2,
\end{align*}
其中$M_{\lambda}=\sup\limits_{n\in\mathbb{Z}}\frac{1}{|(2n\pi)^2-\lambda|}<\infty$。

\end{proof}
\end{example}

\begin{example}
设$\mathscr{X}=C[0,1]$，$A:u(t)\mapsto t\cdot u(t)$，则$A$为有界线性算子，且
\begin{align*}
\sigma(A)=\sigma_r(A)=[0,1].
\end{align*}
\begin{proof}

对任意$\lambda\in[0,1]$，乘法算子$(\lambda-t)^{-1}$为有界线性算子，满足
\begin{align*}
\left\|\frac{1}{\lambda-t}x(t)\right\|\leqslant\sup\limits_{t\in[0,1]}\frac{1}{|\lambda-t|}\|x\|.
\end{align*}
方程$(\lambda-t)u(t)=0$仅有零解。
又$v\in R(\lambda I-A)$需满足$v(\lambda)=0$，故$1\notin\overline{R(\lambda I-A)}$。
因此$[0,1]\subseteq\sigma_r(A)\subseteq\sigma(A)\subseteq[0,1]$，得证。

\end{proof}
\end{example}

\begin{example}
设$\mathscr{X}=L^2[0,1]$，$A:u(t)\mapsto tu(t)$，则$A$为有界线性算子，且
\begin{align*}
\sigma(A)=\sigma_c(A)=[0,1].
\end{align*}
\begin{proof}

与上例类似，$1\notin R(\lambda I-A)$，因$(\lambda-t)^{-1}\notin L^2[0,1]$。
$R(\lambda I-A)$中函数在$t=\lambda$的任意小邻域外可取任意$L^2[0,1]$函数，故$\overline{R(\lambda I-A)}=L^2[0,1]$。

\end{proof}
\end{example}


\subsection{Gelfand定理}

\begin{definition}
设$\mathscr{X}$为$B$空间，$A:D(A)\subseteq\mathscr{X}\to\mathscr{X}$为闭线性算子。\textbf{算子值函数}$R_{\lambda}(A):\rho(A)\to\mathscr{L}(\mathscr{X})$定义为
\begin{align*}
\lambda\mapsto(\lambda I-A)^{-1}\quad(\forall\lambda\in\rho(A)),
\end{align*}
称为$A$的\textbf{预解式}。
\end{definition}

\begin{lemma}\label{lemma:泛函分析--纽曼级数}
设$\mathscr{X}$为$B$空间，$T\in\mathscr{L}(\mathscr{X})$，$\|T\|<1$，则$(I-T)^{-1}\in\mathscr{L}(\mathscr{X})$，且
\begin{align}
\|(I-T)^{-1}\|\leqslant\frac{1}{1-\|T\|}.\label{eq:neuman_est}
\end{align}
还有
\begin{align*}
(I-T)^{-1}=\sum\limits_{k=0}^{\infty}T^k.
\end{align*}
该级数称为$\mathbf{Neuman}$\textbf{级数}.
\end{lemma}
\begin{proof}

对任意固定的 \(y\in\mathscr{X}\)，定义映射 \(S:\mathscr{X}\to\mathscr{X}\) 为
\[
x \mapsto y + Tx,\quad x\in\mathscr{X}.
\]
对任意 \(x,x'\in\mathscr{X}\)，由\refpro{proposition:有界线性算子的基本不等式}有
\[
\|Sx - Sx'\| = \|Tx - Tx'\| \leqslant \|T\|\|x-x'\|.
\]
因为 \(\|T\|<1\)，所以\(S\) 是压缩映射。根据由\refcor{corollary:压缩映射迭代收敛性质}，\(S\) 存在唯一的不动点 \(x\in\mathscr{X}\)，满足 \(x = Sx = y + Tx\)，即 \(x - Tx = y\)，亦即 \((I-T)x = y\)。这表明对每个 \(y\in\mathscr{X}\)，方程 \((I-T)x = y\) 有唯一解 \(x\)，因此 \(I-T\) 是双射，进而 \((I-T)^{-1}y = x\).

由上述证明知,对$\forall y\in \mathscr{X}$,存在唯一的$x\in \mathscr{X}$,使得$(I-T)^{-1}y=x$,从而\(x = y + Tx\)，取范数并利用\refpro{proposition:有界线性算子的基本不等式}可得,对$\forall y\in \mathscr{X}$,有
\[
\|x\| \leqslant \|y\| + \|T\|\|x\|\Longrightarrow (1-\|T\|)\|x\| \leqslant \|y\|\Longrightarrow \|(I-T)^{-1}y\|=\|x\| \leqslant \frac{\|y\|}{1-\|T\|}.
\]
因此
\[
\|(I-T)^{-1}\| = \sup_{\|y\|\leqslant 1}\|(I-T)^{-1}y\| \leqslant \frac{1}{1-\|T\|}.
\]
事实上，由\refcor{corollary:压缩映射迭代收敛性质}知
\begin{align*}
(I-T)^{-1}y=x=\lim\limits_{n\to\infty}S^n y=\sum\limits_{k=0}^{\infty}T^k y,\quad \forall y\in \mathscr{X}.
\end{align*}
即
\begin{align*}
(I-T)^{-1}=\sum\limits_{k=0}^{\infty}T^k.
\end{align*}

\end{proof}

\begin{corollary}
设$\mathscr{X}$为$B$空间，$A$为$\mathscr{X}\to \mathscr{X}$的闭线性算子，则$\rho(A)$是开集,并且对$\forall \lambda_0\in \rho(A)$,当$|\lambda - \lambda_0|<\frac{1}{\|R_{\lambda_0}(A)\|}$时有
\begin{align}\label{eq:rho_open}
R_{\lambda}(A)=\left[ I+\left( \lambda -\lambda _0 \right) R_{\lambda _0}\left( A \right) \right] ^{-1}R_{\lambda _0}\left( A \right) \in \mathscr{L} \left( \mathscr{X} \right) .
\end{align}
\end{corollary}
\begin{proof}

任取$\lambda_0\in\rho(A)$，则
\begin{align*}
\lambda I-A=(\lambda-\lambda_0)I+(\lambda_0 I-A)
=(\lambda_0 I-A)\left[I+(\lambda-\lambda_0)(\lambda_0 I-A)^{-1}\right].
\end{align*}
当$|\lambda-\lambda_0|<\|(\lambda_0 I-A)^{-1}\|^{-1}=\frac{1}{\|R_{\lambda_0}(A)\|}$时，有
\begin{align*}
\|-(\lambda-\lambda_0)(\lambda_0 I-A)^{-1}\|=|\lambda-\lambda_0|\cdot \|(\lambda_0 I-A)^{-1}\|<1,
\end{align*}
则由\reflem{lemma:泛函分析--纽曼级数}知$B\triangleq\left[I+(\lambda-\lambda_0)(\lambda_0 I-A)^{-1}\right]^{-1}\in\mathscr{L}(\mathscr{X})$，则容易验证
\begin{align*}
(\lambda I-A)^{-1}=\left[I+(\lambda-\lambda_0)(\lambda_0 I-A)^{-1}\right]^{-1}(\lambda_0 I-A)^{-1}=BR_{\lambda_0}(A)\in\mathscr{L}(\mathscr{X}),
\end{align*}
故$\lambda\in\rho(A)$,因此$\left( \lambda _0-\frac{1}{\|R_{\lambda_0}(A)\|},\lambda _0+\frac{1}{\|R_{\lambda_0}(A)\|}\right) \subseteq \rho \left( A \right) $,即$\rho(A)$是开集.

\end{proof}

\begin{lemma}[第一预解公式]\label{lemma:第一预解公式}
设$\mathscr{X}$为$B$空间，$A$为$\mathscr{X}\to \mathscr{X}$的闭线性算子，$\lambda,\mu\in\rho(A)$，则
\begin{align}
R_{\lambda}(A)-R_{\mu}(A)=(\mu-\lambda)R_{\lambda}(A)R_{\mu}(A).\label{eq:resolvent_eq}
\end{align}
\end{lemma}
\begin{proof}

直接计算：
\begin{align*}
(\lambda I-A)^{-1}&=(\lambda I-A)^{-1}(\mu I-A)(\mu I-A)^{-1}\\
&=(\lambda I-A)^{-1}\left[(\mu-\lambda)I+\lambda I-A\right](\mu I-A)^{-1}\\
&=(\mu-\lambda)(\lambda I-A)^{-1}(\mu I-A)^{-1}+(\mu I-A)^{-1}.
\end{align*}
整理即得式\eqref{eq:resolvent_eq}。

\end{proof}

\begin{theorem}\label{theorem:预解式都在预解集上是解析的}
设$\mathscr{X}$为$B$空间，$A$为$\mathscr{X}\to \mathscr{X}$的闭线性算子，预解式$R_{\lambda}(A)$在$\rho(A)$内是算子值解析函数。
\end{theorem}
\begin{proof}

(1) 连续性：设$\lambda_0\in\rho(A)$，由式\eqref{eq:rho_open}与\eqref{eq:neuman_est}，当$|\lambda-\lambda_0|<\frac{1}{2\|R_{\lambda_0}(A)\|}$时，
\begin{align*}
\parallel R_{\lambda}(A)\parallel \leqslant \parallel R_{\lambda _0}(A)\parallel \cdot \left\| \left[ I+(\lambda -\lambda _0)R_{\lambda _0}(A) \right] ^{-1} \right\| \leqslant \frac{\parallel R_{\lambda _0}(A)\parallel}{1-\left\| -(\lambda -\lambda _0)R_{\lambda _0}(A) \right\|}\leqslant 2\parallel R_{\lambda _0}(A)\parallel .
\end{align*}
结合\hyperref[lemma:第一预解公式]{第一预解公式}得
\begin{align*}
\|R_{\lambda}(A)-R_{\lambda_0}(A)\|\leqslant|\lambda-\lambda_0|\cdot\|R_{\lambda}(A)\|\cdot\|R_{\lambda_0}(A)\|
\leqslant2\|R_{\lambda_0}(A)\|^2|\lambda-\lambda_0|\to0\quad(\lambda\to\lambda_0).
\end{align*}

(2) 可微性：由\hyperref[lemma:第一预解公式]{第一预解公式}和$R_\lambda(A)$的连续性知
\begin{align*}
\lim\limits_{\lambda\to\lambda_0}\frac{R_{\lambda}(A)-R_{\lambda_0}(A)}{\lambda-\lambda_0}=-\lim\limits_{\lambda\to\lambda_0}R_{\lambda}(A)\cdot R_{\lambda_0}(A)=-\left(R_{\lambda_0}(A)\right)^2.
\end{align*}
故$R_{\lambda}(A)$在$\rho(A)$内解析。

\end{proof}

\begin{theorem}\label{theorem:谱集非空定理}
设$\mathscr{X}$为$B$空间，$A\in \mathscr{L}(\mathscr{X})$，则$\sigma(A)\neq\varnothing $。并且当$|\lambda|>\|A\|$时,有
\begin{align}\label{eq:resolvent_series}
R_{\lambda}(A)=\sum_{n=0}^{\infty}{\frac{1}{\lambda ^{n+1}}A^n},\quad \left\| R_{\lambda}\left( A \right) \right\|\leqslant \frac{1}{\left| \lambda \right|-\left\| A \right\|}.
\end{align}
\end{theorem}
\begin{remark}
因为\hyperref[proposition:有界线性算子必是闭的]{有界线性算子必是闭的},所以$\sigma(A)$是良定义的.
\end{remark}
\begin{proof}

反证法：若$\rho(A)=\mathbb{C}$，则由\refthe{theorem:预解式都在预解集上是解析的}知$R_{\lambda}(A)$在$\mathbb{C}$上解析。
当$|\lambda|>\|A\|$时，有$\|\frac{A}{\lambda}\|=\frac{\|A\|}{|\lambda|}<1$.从而由\reflem{lemma:泛函分析--纽曼级数}知其Neuman级数展开为
\begin{align*}
R_{\lambda}(A)=\left( \lambda I-A \right) ^{-1}=\frac{1}{\lambda} \left( I-\frac{A}{\lambda} \right) ^{-1}=\sum_{n=0}^{\infty}{\frac{1}{\lambda ^{n+1}}A^n},
\end{align*}
且
\begin{align*}
\left\| R_{\lambda}\left( A \right) \right\| =\frac{1}{\left| \lambda \right|}\left\| \left( I-\frac{A}{\lambda} \right) ^{-1} \right\| \leqslant \frac{1}{\left| \lambda \right|}\frac{1}{1-\left\| \frac{A}{\lambda} \right\|}=\frac{1}{\left| \lambda \right|-\left\| A \right\|}.
\end{align*}
故$\|R_{\lambda}(A)\|$在复平面上有界，不妨设$\|R_{\lambda}(A)\|\leqslant M,\forall \lambda \in \mathbb{C}$,其中$M>0$.

任取$f\in\mathscr{L}(\mathscr{X})^*$，定义数值解析函数
\begin{align*}
u_f(\lambda)\triangleq f(R_{\lambda}(A)),\,\,\forall \lambda \in \mathbb{C}.
\end{align*}
从而由\refpro{proposition:有界线性算子的基本不等式}知
\begin{align*}
\left| u_f\left( \lambda \right) \right|=\left| f\left( R_{\lambda}\left( A \right) \right) \right|\leqslant \left\| f \right\| \cdot \left\| R_{\lambda}\left( A \right) \right\| \leqslant M\left\| f \right\| ,\quad \forall \lambda \in \mathbb{C} .
\end{align*}
由\hyperref[Complex-Analysis-theorem:Liouville(刘维尔)定理]{Liouville定理}，$u_f(\lambda)$为常值函数，故$R_{\lambda}(A)$为常值算子，与\hyperref[lemma:第一预解公式]{第一预解公式}矛盾!

\end{proof}

\begin{definition}[谱半径]
设$\mathscr{X}$为$B$空间,$A\in\mathscr{L}(\mathscr{X})$，称数
\begin{align*}
r_{\sigma}(A)\triangleq\sup\{|\lambda|\mid\lambda\in\sigma(A)\}
\end{align*}
为$A$的\textbf{谱半径}。
\end{definition}

\begin{proposition}\label{proposition:算子范数都大于谱半径}
设$\mathscr{X}$为$B$空间,$A\in\mathscr{L}(\mathscr{X})$,则$r_{\sigma}(A)\leqslant\|A\|$.
\end{proposition}
\begin{proof}

对$\forall \lambda \in \mathbb{C}$，若$|\lambda|>\|A\|$，则$\left\| \frac{A}{\lambda} \right\| = \frac{\|A\|}{|\lambda|}<1$。由引理知$\left( I-\frac{A}{\lambda} \right)^{-1} \in \mathscr{L}(\mathscr{X})$。于是容易发现$(\lambda I - A)^{-1} = \frac{1}{\lambda}\left( I-\frac{A}{\lambda} \right)^{-1} \in \mathscr{L}(\mathscr{X})$，即$\lambda \in \rho(A)$，故$\lambda \notin \sigma(A)$。因此$|\lambda| \leqslant \|A\|$，$\forall \lambda \in \sigma(A)$。进而$r_{\sigma}(A) \leqslant \|A\|$。
\end{proof}

\begin{theorem}[Gelfand定理]\label{theorem:泛函分析--Gelfand定理}
设$\mathscr{X}$是$B$空间，$A\in\mathscr{L}(\mathscr{X})$，则
\begin{align}
r_{\sigma}(A)=\lim\limits_{n\to\infty}\|A^n\|^{\frac{1}{n}}.\label{eq:gelfand_formula}
\end{align}
\end{theorem}
\begin{remark}
上述定理结论可推广到一般Banach代数。
\end{remark}
\begin{proof}

首先，由\eqref{eq:resolvent_series}式与\hyperref[Complex-Analysis-theorem:Cauchy-Hadamard收敛半径公式]{Cauchy-Hadamard收敛半径公式}，当$|\lambda|>\varlimsup\limits_{n\to\infty}\|A^n\|^{\frac{1}{n}}$时，$R_{\lambda}(A)\in\mathscr{L}(\mathscr{X})$，故再由\refpro{proposition:算子范数都大于谱半径}知
\begin{align}
r_{\sigma}(A)\leqslant\varlimsup\limits_{n\to\infty}\|A^n\|^{\frac{1}{n}}.\label{eq:upper_bound}
\end{align}

记$a\triangleq r_{\sigma}(A)$，任取$f\in\mathscr{L}(\mathscr{X})^*$，定义$u_f(\lambda)\triangleq f(R_{\lambda}(A))$，其在$|\lambda|>a$解析，且有Laurent展开
\begin{align*}
u_f(\lambda)=\sum\limits_{n=0}^{\infty}\frac{1}{\lambda^{n+1}}f(A^n).
\end{align*}
由Laurent展开收敛半径性质，对任意$\varepsilon>0$，
\begin{align*}
\sum\limits_{n=0}^{\infty}\frac{1}{(a+\varepsilon)^{n+1}}|f(A^n)|<\infty,
\end{align*}
即$\left|f\left(\frac{A^n}{(a+\varepsilon)^{n+1}}\right)\right|$对任意$f\in\mathscr{L}(\mathscr{X})^*$有界。
由\hyperref[theorem:泛函分析--共鸣定理]{共鸣定理}，存在常数$M>0$，使得
\begin{align*}
\frac{1}{(a+\varepsilon)^{n+1}}\|A^n\|\leqslant M,
\end{align*}
故$\varlimsup\limits_{n\to\infty}\|A^n\|^{\frac{1}{n}}\leqslant a+\varepsilon$。
由$\varepsilon$的任意性，
\begin{align}
\varlimsup\limits_{n\to\infty}\|A^n\|^{\frac{1}{n}}\leqslant a.\label{eq:lower_bound}
\end{align}
结合\eqref{eq:upper_bound}与\eqref{eq:lower_bound}得$r_{\sigma}(A)=\varlimsup\limits_{n\to\infty}\|A^n\|^{\frac{1}{n}}$。

另一方面，对任意$\lambda\in\mathbb{C}$，$\lambda^n I-A^n=(\lambda I-A)P_{\lambda}(A)=P_{\lambda}(A)(\lambda I-A)$，其中$P_{\lambda}(A)=\sum\limits_{j=1}^{n}\lambda^{j-1}A^{n-j}$。
若$\lambda^n\in\rho(A^n)$，则$\lambda\in\rho(A)$，故$\lambda\in\sigma(A)\implies\lambda^n\in\sigma(A^n)$，即$|\lambda|^n\leqslant\|A^n\|$，得
\begin{align}
r_{\sigma}(A)\leqslant\lim\limits_{n\to\infty}\|A^n\|^{\frac{1}{n}}.\label{eq:lim_bound}
\end{align}
因此$\lim\limits_{n\to\infty}\|A^n\|^{\frac{1}{n}}$存在，且等于$r_{\sigma}(A)$。

\end{proof}


\subsection{例子}

\begin{example}
在空间$l^2$上定义右推移算子$A$：
\begin{align*}
A:x=(x_1,x_2,\cdots,x_n,\cdots)\mapsto(0,x_1,x_2,\cdots,x_n,\cdots).
\end{align*}
则该算子的谱满足：
\begin{align*}
\sigma_p(A)=\varnothing ,\quad \sigma_r(A)=\{\lambda\in\mathbb{C}\mid|\lambda|<1\},\quad \sigma_c(A)=\{\lambda\in\mathbb{C}\mid|\lambda|=1\}.
\end{align*}
\end{example}
\begin{proof}

由$\|A\|=1$，结合Gelfand定理可得
\begin{align*}
\sigma(A)\subseteq\{\lambda\in\mathbb{C}\mid|\lambda|\leqslant1\}.
\end{align*}
可验证$A$的共轭算子为$A^*:l^2\to l^2$，且
\begin{align*}
A^*x=(x_2,x_3,\cdots,x_{n+1},\cdots).
\end{align*}
对任意$\lambda\in\mathbb{C}$，有
\begin{align*}
((\lambda I-A)x,y)=(x,(\overline{\lambda}I-A^*)y),\quad \forall x,y\in l^2,
\end{align*}
由此推得
\begin{align*}
\overline{R(\lambda I-A)}^{\perp}=N(\overline{\lambda}I-A^*).
\end{align*}
\begin{enumerate}[(1)]
\item 证明$\sigma_p(A)=\varnothing $。
设$(\lambda I-A)x=0$，则$\lambda x_n=x_{n-1},\ n\geqslant2$，$\lambda x_1=0$。
- 若$\lambda\neq0$，则$x=0$；
- 若$\lambda=0$，则$x=0$。
因此$(\lambda I-A)^{-1}$恒存在，故$\sigma_p(A)=\varnothing $。

\item 证明$\sigma_r(A)=\{\lambda\in\mathbb{C}\mid|\lambda|<1\}$，$\sigma_c(A)=\{\lambda\in\mathbb{C}\mid|\lambda|=1\}$。
只需验证：
\begin{align*}
&|\lambda|<1\implies \overline{R(\lambda I-A)}^{\perp}=N(\overline{\lambda}I-A^*)\neq\{0\},\\
&|\lambda|=1\implies \overline{R(\lambda I-A)}^{\perp}=N(\overline{\lambda}I-A^*)=\{0\}.
\end{align*}
求解方程$(\overline{\lambda}I-A^*)x=0$，得$\overline{\lambda}x_n=x_{n+1}$，即$x_{n+1}=\overline{\lambda}^n x_1,\ n=1,2,\cdots$。
- 当$|\lambda|<1$时，$N(\overline{\lambda}I-A^*)=\{c(1,\overline{\lambda},\overline{\lambda}^2,\cdots,\overline{\lambda}^n,\cdots)\mid c\in\mathbb{C}\}$，故$N(\overline{\lambda}I-A^*)\neq\{0\}$，$R(\lambda I-A)$在$l^2$中不稠密，即$\lambda\in\sigma_r(A)$，因此$\{\lambda\in\mathbb{C}\mid|\lambda|<1\}\subseteq\sigma_r(A)$。结合$\sigma(A)=\{\lambda\in\mathbb{C}\mid|\lambda|\leqslant1\}$得该部分结论。
- 当$|\lambda|=1$时，$x=x_1(1,\overline{\lambda},\cdots,\overline{\lambda}^n,\cdots)\notin l^2$，故$N(\overline{\lambda}I-A^*)=\{0\}$，即$R(\lambda I-A)$在$l^2$中稠密。结合$\lambda\in\sigma(A)$，得$\lambda\in\sigma_c(A)$。
\end{enumerate}

\end{proof}

\begin{definition}
设$(\mathscr{X},(\cdot,\cdot))$为Hilbert空间，算子$A\in\mathscr{L}(\mathscr{X})$，若满足$A^*=A$，即
\begin{align}
(Ax,y)=(x,Ay),\quad \forall x,y\in\mathscr{X}.
\label{eq:sym-op}
\end{align}
则称$A$为\textbf{对称算子}。
\end{definition}

\begin{definition}
设$(\mathscr{X},(\cdot,\cdot))$为Hilbert空间，算子$U\in\mathscr{L}(\mathscr{X})$，若满足
\begin{align}
(Ux,Uy)=(x,y),\ \forall x,y\in\mathscr{X},\quad R(U)=\mathscr{X}.
\label{eq:unitary-op}
\end{align}
则称$U$为\textbf{酉算子}。
\end{definition}

\begin{proposition}
设$(\mathscr{X},(\cdot,\cdot))$为Hilbert空间，$A\in\mathscr{L}(\mathscr{X})$为对称算子，则$\sigma(A)\subseteq\mathbb{R}$，且$\sigma_r(A)=\varnothing $。
\end{proposition}
\begin{proof}

由对称算子定义，$(Ax,x)\in\mathbb{R},\ \forall x\in\mathscr{X}$。设$\lambda=a+ib,\ b\neq0$，先证$\lambda I-A$存在有界逆。
\begin{align}
\|(\lambda I-A)x\|^2=\|(aI-A)x\|^2+b^2\|x\|^2.
\label{eq:sym-norm}
\end{align}
因此$N(\lambda I-A)=\{0\}$，即$\lambda I-A$为单射。

下证$\lambda I-A$为满射，即$R(\lambda I-A)=\mathscr{X}$。
首先证明$R(\lambda I-A)$是闭集。设$\{x_n\}\subseteq\mathscr{X}$，满足$(\lambda I-A)x_n=y_n\to y\ (n\to\infty)$，由式\eqref{eq:sym-norm}得
\begin{align*}
\|(\lambda I-A)(x_n-x_m)\|^2\geqslant b^2\|x_n-x_m\|^2,
\end{align*}
故$\{x_n\}$为$\mathscr{X}$中Cauchy列，设$x_n\to x\ (n\to\infty)$，则
\begin{align*}
y_n=(\lambda I-A)x_n\to(\lambda I-A)x=y,
\end{align*}
因此$R(\lambda I-A)=\overline{R(\lambda I-A)}$，即值域闭。

再证$R(\lambda I-A)^{\perp}=\{0\}$。设$y\in R(\lambda I-A)^{\perp}$，则$((\lambda I-A)x,y)=0,\ \forall x\in\mathscr{X}$。
由$A$对称，得$(x,(\overline{\lambda}I-A)y)=0,\ \forall x\in\mathscr{X}$，故$(\overline{\lambda}I-A)y=0$，即$y\in N(\overline{\lambda}I-A)$。
将式\eqref{eq:sym-norm}中$\lambda$替换为$\overline{\lambda}$，得$N(\overline{\lambda}I-A)=\{0\}$，故$y=0$，即$R(\lambda I-A)^{\perp}=\{0\}$。

综上$R(\lambda I-A)=\mathscr{X}$，由\hyperref[theorem:泛函分析--Banach逆算子定理]{Banach逆算子定理}，$\lambda I-A$存在有界逆，即$\lambda\notin\sigma(A)$，故$\sigma(A)\subseteq\mathbb{R}$。

设$\lambda\in\sigma(A)$且$\lambda\notin\sigma_p(A)$，则$N(\lambda I-A)=\{0\}$，此时
\begin{align*}
R(\lambda I-A)^{\perp}=N(\overline{\lambda}I-A^*)=N(\lambda I-A)=\{0\},
\end{align*}
即$R(\lambda I-A)$在$\mathscr{X}$中稠密，故$\lambda\in\sigma_c(A)$，因此$\sigma_r(A)=\varnothing $。

\end{proof}

\begin{proposition}
设$(\mathscr{X},(\cdot,\cdot))$为Hilbert空间，$U\in\mathscr{L}(\mathscr{X})$为酉算子，则$\sigma(U)\subseteq S^1=\{e^{i\theta}\mid\theta\in[0,2\pi]\}$，且$\sigma_r(U)=\varnothing $。
\end{proposition}
\begin{proof}

由酉算子定义，$\|Ux\|^2=\|x\|^2$，故$U$为单射，且$\|U\|=1$，因此$\sigma(U)\subseteq\{\lambda\in\mathbb{C}\mid|\lambda|\leqslant1\}$。
又$U$满射，由\hyperref[theorem:泛函分析--Banach逆算子定理]{Banach逆算子定理}，$U$存在有界逆$U^{-1}$且$\|U^{-1}\|=1$。

若$|\lambda|<1$，则
\begin{align}
\lambda I-U=-U(I-\lambda U^{-1}).
\label{eq:unitary-decomp}
\end{align}
因$\|\lambda U^{-1}\|=|\lambda|\cdot\|U^{-1}\|=|\lambda|<1$，故$I-\lambda U^{-1}$存在有界逆，因此$\lambda I-U$存在有界逆，且
\begin{align*}
(\lambda I-U)^{-1}=(I-\lambda U^{-1})^{-1}U^{-1},
\end{align*}
即$\lambda\notin\sigma(U)$，故$\sigma(U)\subseteq\{\lambda\in\mathbb{C}\mid|\lambda|=1\}$。

设$\lambda\in\sigma(U)$且$\lambda\notin\sigma_p(U)$，则$N(\lambda I-U)=\{0\}$，此时
\begin{align*}
R(\lambda I-U)^{\perp}=N(\overline{\lambda}I-U^{-1})=N(\lambda I-U)=\{0\},
\end{align*}
即$R(\lambda I-U)$在$\mathscr{X}$中稠密，故$\lambda\in\sigma_c(U)$，因此$\sigma_r(U)=\varnothing $。

\end{proof}



\begin{example}
设$\mathscr{X}$是$B$空间，求证:$\mathscr{L}(\mathscr{X})$中可逆(存在有界逆)算子集是开的。
\end{example}

\begin{example}
设$A$是闭线性算子，$\lambda_1,\lambda_2,\cdots,\lambda_n\in\sigma_p(A)$两两互异，又设$x_i$是对应于$\lambda_i$的特征元$(i=1,2,\cdots,n)$。求证:$\{x_1,x_2,\cdots,x_n\}$是线性无关的。
\end{example}

\begin{example}
在双边$l^2$空间上，考察右推移算子
\begin{align*}
A:x=(\cdots,\xi_{-n},\xi_{-n+1},\cdots,\xi_{-1},\xi_0,\xi_1,\cdots,\xi_{n-1},\xi_n,\cdots)\in l^2
\mapsto Ax=(\cdots,\eta_{-n},\eta_{-n+1},\cdots,\eta_{-1},\eta_0,\eta_1,\cdots,\eta_{n-1},\eta_n,\cdots),
\end{align*}
其中$\eta_m=\xi_{m-1}(m\in\mathbb{Z})$。求证:$\sigma_c(A)=\sigma(A)=$单位圆周。
\end{example}

\begin{example}
在$l^2$空间上，考察左推移算子
\begin{align*}
A:(\xi_1,\xi_2,\cdots)\mapsto(\xi_2,\xi_3,\cdots).
\end{align*}
求证:$\sigma_p(A)=\{\lambda\in\mathbb{C}\mid|\lambda|<1\},\sigma_c(A)=\{\lambda\in\mathbb{C}\mid|\lambda|=1\}$，并且$\sigma(A)=\sigma_p(A)\cup\sigma_c(A)$。
\end{example}




\end{document}